% ==============================================================================
% Centralised Baseline section -- Overleaf-ready, drop-in.
% Reporting style: descriptive (mean +/- SD across seeds), bolded best where
% comparing methods, no formal hypothesis testing -- matches the FL literature
% convention used throughout this dissertation.
%
% Required tables (none external -- table is inline in this .tex).
% Required figures: none for this section.
%
% Cross-references used:
%   sec:dataset       -- DermaMNIST dataset description in Methods chapter
%   sec:model         -- DermMNISTCNN architecture in Methods chapter
%   sec:metrics       -- metric definitions in Methods chapter
%   sec:fl-results    -- federated-learning results chapter (forward reference)
%
% Usage: % ==============================================================================
% Centralised Baseline section -- Overleaf-ready, drop-in.
% Reporting style: descriptive (mean +/- SD across seeds), bolded best where
% comparing methods, no formal hypothesis testing -- matches the FL literature
% convention used throughout this dissertation.
%
% Required tables (none external -- table is inline in this .tex).
% Required figures: none for this section.
%
% Cross-references used:
%   sec:dataset       -- DermaMNIST dataset description in Methods chapter
%   sec:model         -- DermMNISTCNN architecture in Methods chapter
%   sec:metrics       -- metric definitions in Methods chapter
%   sec:fl-results    -- federated-learning results chapter (forward reference)
%
% Usage: % ==============================================================================
% Centralised Baseline section -- Overleaf-ready, drop-in.
% Reporting style: descriptive (mean +/- SD across seeds), bolded best where
% comparing methods, no formal hypothesis testing -- matches the FL literature
% convention used throughout this dissertation.
%
% Required tables (none external -- table is inline in this .tex).
% Required figures: none for this section.
%
% Cross-references used:
%   sec:dataset       -- DermaMNIST dataset description in Methods chapter
%   sec:model         -- DermMNISTCNN architecture in Methods chapter
%   sec:metrics       -- metric definitions in Methods chapter
%   sec:fl-results    -- federated-learning results chapter (forward reference)
%
% Usage: % ==============================================================================
% Centralised Baseline section -- Overleaf-ready, drop-in.
% Reporting style: descriptive (mean +/- SD across seeds), bolded best where
% comparing methods, no formal hypothesis testing -- matches the FL literature
% convention used throughout this dissertation.
%
% Required tables (none external -- table is inline in this .tex).
% Required figures: none for this section.
%
% Cross-references used:
%   sec:dataset       -- DermaMNIST dataset description in Methods chapter
%   sec:model         -- DermMNISTCNN architecture in Methods chapter
%   sec:metrics       -- metric definitions in Methods chapter
%   sec:fl-results    -- federated-learning results chapter (forward reference)
%
% Usage: \input{15_centralised_baseline}
% ==============================================================================

\section{Centralised Baseline}
\label{sec:centralised-baseline}

Before evaluating any federated configuration, we establish a
centralised non-federated reference. The centralised baseline trains
the same \texttt{DermMNISTCNN} architecture
(Section~\ref{sec:model}) on the pooled DermaMNIST training set with
no federation. It provides an upper bound on what can be achieved on
this task with this architecture, against which the federated
results in Section~\ref{sec:fl-results} are interpreted as
fractional recoveries of the centralised performance.

\subsection{Training setup}
\label{sec:centralised-setup}

The centralised baseline pools all 7{,}007 training samples, the
1{,}003 validation samples, and the 2{,}005 test samples described
in Section~\ref{sec:dataset} into a single non-federated training
loop. The model architecture is held identical to the federated
experiments (Section~\ref{sec:model}): a four-block GroupNorm CNN
with 423{,}175 trainable parameters operating on $28\times28$ RGB
inputs. The optimiser is SGD with learning rate $0.01$, momentum
$0.9$, weight decay $0.0$, and batch size $32$ --- the same
hyperparameters used by each client in the federated experiments,
chosen for direct comparability rather than for centralised
optimum. The model is trained for $50$ epochs, with the test set
evaluated once at the epoch of highest validation macro-F1
(\emph{best-validation checkpoint} protocol).

To match the federated experiments and support a paired-seed
analysis where applicable, the centralised baseline is repeated
across the same ten random seeds used for the federated runs:
$\{42,\,123,\,456,\,789,\,999,\,2024,\,31337,\,8675309,\,161803,\,$
$271828\}$. Within each seed the data shuffle, model initialisation,
and training trajectory are deterministic; the seed-level variation
quantifies optimisation noise on this task.

\subsection{Results}
\label{sec:centralised-results}

Table~\ref{tab:centralised-results} reports test performance at the
best-validation checkpoint across the ten seeds. The centralised
model reaches a mean test macro-F1 of $0.560\pm0.024$ with mean
classification accuracy $0.764\pm0.012$ and mean balanced accuracy
$0.561\pm0.034$. The per-seed test macro-F1 ranges from $0.518$ to
$0.597$. The best-validation epoch is consistently in the second
half of training, with a median of $47$ out of $50$ epochs, so the
$50$-epoch training budget is appropriate for this task and
architecture.

\begin{table}[!htbp]
\centering
\caption{Centralised baseline test performance at the
best-validation checkpoint, reported as mean $\pm$ standard
deviation across $n=10$ seeds. Per-class F1 is the F1 score on
each of the seven DermaMNIST classes; class prevalence in the
training set is shown in parentheses.}
\label{tab:centralised-results}
\small
\begin{tabular}{lc}
\toprule
Metric & Test value (mean $\pm$ SD) \\
\midrule
Macro-F1                        & $\mathbf{0.5600 \pm 0.0240}$ \\
Balanced accuracy               & $0.5608 \pm 0.0341$ \\
Overall accuracy                & $0.7637 \pm 0.0123$ \\
Cross-entropy loss              & $0.7778 \pm 0.0653$ \\
\midrule
\multicolumn{2}{l}{\textit{Per-class F1:}} \\
\quad Actinic keratoses (3.27\%)              & $0.4923 \pm 0.0433$ \\
\quad Basal cell carcinoma (5.13\%)           & $0.5380 \pm 0.0263$ \\
\quad Benign keratosis (10.97\%)              & $0.5285 \pm 0.0397$ \\
\quad Dermatofibroma (1.15\%)                 & $0.3359 \pm 0.0961$ \\
\quad Melanoma (11.11\%)                      & $0.4551 \pm 0.0349$ \\
\quad Melanocytic nevi (67.05\%)              & $0.8896 \pm 0.0088$ \\
\quad Vascular lesions (1.41\%)               & $0.6807 \pm 0.0680$ \\
\bottomrule
\end{tabular}
\end{table}

The per-class breakdown reveals the task's structure clearly. The
majority class \emph{melanocytic nevi} (67\% prevalence) is the
easiest to learn, reaching $F1 = 0.890$ with very small across-seed
variance ($\pm 0.009$). The two rarest classes,
\emph{dermatofibroma} ($1.15\%$) and \emph{vascular lesions}
($1.41\%$), show the largest across-seed variance ($\pm 0.096$ and
$\pm 0.068$), reflecting the difficulty of stably learning classes
represented by only $80$ and $99$ training samples respectively.
The two clinically critical classes \emph{melanoma} and
\emph{basal cell carcinoma} reach mean F1 of $0.455$ and $0.538$
respectively --- well above chance but substantially below the
majority class, consistent with the medical-imaging literature on
small CNNs at $28\times28$ resolution.

The macro-F1 of $0.56$ is a benchmark-typical result for
DermaMNIST at this resolution and architecture, comparable to
published values on the MedMNIST v2 benchmark. It is not intended
as a clinical-grade performance estimate; the input resolution is
deliberately compressed and the architecture is a small CNN, both
chosen for comparability across the federated experiments in this
dissertation rather than for clinical deployment.

\subsection{Role in the subsequent analysis}
\label{sec:centralised-role}

The centralised baseline serves three purposes throughout the rest
of the dissertation.

First, it is the \emph{upper bound} against which federated
configurations are measured. Subsequent results report recovery
percentages relative to the $0.5600$ centralised mean
(e.g.\ ``FedProx recovers $90.7\%$ of centralised performance''),
which makes the absolute magnitude of federation costs
interpretable. The within-federation algorithmic differences
discussed in Sections~\ref{sec:results-statistical-het}
and~\ref{sec:results-system-het} are then framed against this
fixed reference rather than against each other in isolation.

Second, the across-seed variance of the centralised baseline
($\pm 0.024$ macro-F1) sets a noise floor below which seed-level
differences in federated experiments cannot be reliably
distinguished from optimisation noise alone. Federated results
with within-pair $\Delta$ values smaller than this in magnitude
should be read as descriptive rather than substantive.

Third, the per-class profile in
Table~\ref{tab:centralised-results} establishes the
\emph{intrinsic difficulty} of each class on this task,
independent of any federation effect. The federated per-class
breakdowns reported later in the dissertation can be interpreted
relative to this profile: a class that already shows high
variance in the centralised setting is unlikely to support strong
inferential claims in the federated setting either.

\subsection{Summary}
\label{sec:centralised-summary}

The centralised baseline establishes that the $28\times28$
DermaMNIST classification task with the \texttt{DermMNISTCNN}
architecture supports a mean test macro-F1 of approximately $0.56$
under non-federated training, with the majority class melanocytic
nevi at near-ceiling performance ($F1 \approx 0.89$) and the
rarest classes at substantially lower and noisier per-class F1.
This serves as the reference against which all federated results
in the remainder of the dissertation are interpreted.

% ==============================================================================

\section{Centralised Baseline}
\label{sec:centralised-baseline}

Before evaluating any federated configuration, we establish a
centralised non-federated reference. The centralised baseline trains
the same \texttt{DermMNISTCNN} architecture
(Section~\ref{sec:model}) on the pooled DermaMNIST training set with
no federation. It provides an upper bound on what can be achieved on
this task with this architecture, against which the federated
results in Section~\ref{sec:fl-results} are interpreted as
fractional recoveries of the centralised performance.

\subsection{Training setup}
\label{sec:centralised-setup}

The centralised baseline pools all 7{,}007 training samples, the
1{,}003 validation samples, and the 2{,}005 test samples described
in Section~\ref{sec:dataset} into a single non-federated training
loop. The model architecture is held identical to the federated
experiments (Section~\ref{sec:model}): a four-block GroupNorm CNN
with 423{,}175 trainable parameters operating on $28\times28$ RGB
inputs. The optimiser is SGD with learning rate $0.01$, momentum
$0.9$, weight decay $0.0$, and batch size $32$ --- the same
hyperparameters used by each client in the federated experiments,
chosen for direct comparability rather than for centralised
optimum. The model is trained for $50$ epochs, with the test set
evaluated once at the epoch of highest validation macro-F1
(\emph{best-validation checkpoint} protocol).

To match the federated experiments and support a paired-seed
analysis where applicable, the centralised baseline is repeated
across the same ten random seeds used for the federated runs:
$\{42,\,123,\,456,\,789,\,999,\,2024,\,31337,\,8675309,\,161803,\,$
$271828\}$. Within each seed the data shuffle, model initialisation,
and training trajectory are deterministic; the seed-level variation
quantifies optimisation noise on this task.

\subsection{Results}
\label{sec:centralised-results}

Table~\ref{tab:centralised-results} reports test performance at the
best-validation checkpoint across the ten seeds. The centralised
model reaches a mean test macro-F1 of $0.560\pm0.024$ with mean
classification accuracy $0.764\pm0.012$ and mean balanced accuracy
$0.561\pm0.034$. The per-seed test macro-F1 ranges from $0.518$ to
$0.597$. The best-validation epoch is consistently in the second
half of training, with a median of $47$ out of $50$ epochs, so the
$50$-epoch training budget is appropriate for this task and
architecture.

\begin{table}[!htbp]
\centering
\caption{Centralised baseline test performance at the
best-validation checkpoint, reported as mean $\pm$ standard
deviation across $n=10$ seeds. Per-class F1 is the F1 score on
each of the seven DermaMNIST classes; class prevalence in the
training set is shown in parentheses.}
\label{tab:centralised-results}
\small
\begin{tabular}{lc}
\toprule
Metric & Test value (mean $\pm$ SD) \\
\midrule
Macro-F1                        & $\mathbf{0.5600 \pm 0.0240}$ \\
Balanced accuracy               & $0.5608 \pm 0.0341$ \\
Overall accuracy                & $0.7637 \pm 0.0123$ \\
Cross-entropy loss              & $0.7778 \pm 0.0653$ \\
\midrule
\multicolumn{2}{l}{\textit{Per-class F1:}} \\
\quad Actinic keratoses (3.27\%)              & $0.4923 \pm 0.0433$ \\
\quad Basal cell carcinoma (5.13\%)           & $0.5380 \pm 0.0263$ \\
\quad Benign keratosis (10.97\%)              & $0.5285 \pm 0.0397$ \\
\quad Dermatofibroma (1.15\%)                 & $0.3359 \pm 0.0961$ \\
\quad Melanoma (11.11\%)                      & $0.4551 \pm 0.0349$ \\
\quad Melanocytic nevi (67.05\%)              & $0.8896 \pm 0.0088$ \\
\quad Vascular lesions (1.41\%)               & $0.6807 \pm 0.0680$ \\
\bottomrule
\end{tabular}
\end{table}

The per-class breakdown reveals the task's structure clearly. The
majority class \emph{melanocytic nevi} (67\% prevalence) is the
easiest to learn, reaching $F1 = 0.890$ with very small across-seed
variance ($\pm 0.009$). The two rarest classes,
\emph{dermatofibroma} ($1.15\%$) and \emph{vascular lesions}
($1.41\%$), show the largest across-seed variance ($\pm 0.096$ and
$\pm 0.068$), reflecting the difficulty of stably learning classes
represented by only $80$ and $99$ training samples respectively.
The two clinically critical classes \emph{melanoma} and
\emph{basal cell carcinoma} reach mean F1 of $0.455$ and $0.538$
respectively --- well above chance but substantially below the
majority class, consistent with the medical-imaging literature on
small CNNs at $28\times28$ resolution.

The macro-F1 of $0.56$ is a benchmark-typical result for
DermaMNIST at this resolution and architecture, comparable to
published values on the MedMNIST v2 benchmark. It is not intended
as a clinical-grade performance estimate; the input resolution is
deliberately compressed and the architecture is a small CNN, both
chosen for comparability across the federated experiments in this
dissertation rather than for clinical deployment.

\subsection{Role in the subsequent analysis}
\label{sec:centralised-role}

The centralised baseline serves three purposes throughout the rest
of the dissertation.

First, it is the \emph{upper bound} against which federated
configurations are measured. Subsequent results report recovery
percentages relative to the $0.5600$ centralised mean
(e.g.\ ``FedProx recovers $90.7\%$ of centralised performance''),
which makes the absolute magnitude of federation costs
interpretable. The within-federation algorithmic differences
discussed in Sections~\ref{sec:results-statistical-het}
and~\ref{sec:results-system-het} are then framed against this
fixed reference rather than against each other in isolation.

Second, the across-seed variance of the centralised baseline
($\pm 0.024$ macro-F1) sets a noise floor below which seed-level
differences in federated experiments cannot be reliably
distinguished from optimisation noise alone. Federated results
with within-pair $\Delta$ values smaller than this in magnitude
should be read as descriptive rather than substantive.

Third, the per-class profile in
Table~\ref{tab:centralised-results} establishes the
\emph{intrinsic difficulty} of each class on this task,
independent of any federation effect. The federated per-class
breakdowns reported later in the dissertation can be interpreted
relative to this profile: a class that already shows high
variance in the centralised setting is unlikely to support strong
inferential claims in the federated setting either.

\subsection{Summary}
\label{sec:centralised-summary}

The centralised baseline establishes that the $28\times28$
DermaMNIST classification task with the \texttt{DermMNISTCNN}
architecture supports a mean test macro-F1 of approximately $0.56$
under non-federated training, with the majority class melanocytic
nevi at near-ceiling performance ($F1 \approx 0.89$) and the
rarest classes at substantially lower and noisier per-class F1.
This serves as the reference against which all federated results
in the remainder of the dissertation are interpreted.

% ==============================================================================

\section{Centralised Baseline}
\label{sec:centralised-baseline}

Before evaluating any federated configuration, we establish a
centralised non-federated reference. The centralised baseline trains
the same \texttt{DermMNISTCNN} architecture
(Section~\ref{sec:model}) on the pooled DermaMNIST training set with
no federation. It provides an upper bound on what can be achieved on
this task with this architecture, against which the federated
results in Section~\ref{sec:fl-results} are interpreted as
fractional recoveries of the centralised performance.

\subsection{Training setup}
\label{sec:centralised-setup}

The centralised baseline pools all 7{,}007 training samples, the
1{,}003 validation samples, and the 2{,}005 test samples described
in Section~\ref{sec:dataset} into a single non-federated training
loop. The model architecture is held identical to the federated
experiments (Section~\ref{sec:model}): a four-block GroupNorm CNN
with 423{,}175 trainable parameters operating on $28\times28$ RGB
inputs. The optimiser is SGD with learning rate $0.01$, momentum
$0.9$, weight decay $0.0$, and batch size $32$ --- the same
hyperparameters used by each client in the federated experiments,
chosen for direct comparability rather than for centralised
optimum. The model is trained for $50$ epochs, with the test set
evaluated once at the epoch of highest validation macro-F1
(\emph{best-validation checkpoint} protocol).

To match the federated experiments and support a paired-seed
analysis where applicable, the centralised baseline is repeated
across the same ten random seeds used for the federated runs:
$\{42,\,123,\,456,\,789,\,999,\,2024,\,31337,\,8675309,\,161803,\,$
$271828\}$. Within each seed the data shuffle, model initialisation,
and training trajectory are deterministic; the seed-level variation
quantifies optimisation noise on this task.

\subsection{Results}
\label{sec:centralised-results}

Table~\ref{tab:centralised-results} reports test performance at the
best-validation checkpoint across the ten seeds. The centralised
model reaches a mean test macro-F1 of $0.560\pm0.024$ with mean
classification accuracy $0.764\pm0.012$ and mean balanced accuracy
$0.561\pm0.034$. The per-seed test macro-F1 ranges from $0.518$ to
$0.597$. The best-validation epoch is consistently in the second
half of training, with a median of $47$ out of $50$ epochs, so the
$50$-epoch training budget is appropriate for this task and
architecture.

\begin{table}[!htbp]
\centering
\caption{Centralised baseline test performance at the
best-validation checkpoint, reported as mean $\pm$ standard
deviation across $n=10$ seeds. Per-class F1 is the F1 score on
each of the seven DermaMNIST classes; class prevalence in the
training set is shown in parentheses.}
\label{tab:centralised-results}
\small
\begin{tabular}{lc}
\toprule
Metric & Test value (mean $\pm$ SD) \\
\midrule
Macro-F1                        & $\mathbf{0.5600 \pm 0.0240}$ \\
Balanced accuracy               & $0.5608 \pm 0.0341$ \\
Overall accuracy                & $0.7637 \pm 0.0123$ \\
Cross-entropy loss              & $0.7778 \pm 0.0653$ \\
\midrule
\multicolumn{2}{l}{\textit{Per-class F1:}} \\
\quad Actinic keratoses (3.27\%)              & $0.4923 \pm 0.0433$ \\
\quad Basal cell carcinoma (5.13\%)           & $0.5380 \pm 0.0263$ \\
\quad Benign keratosis (10.97\%)              & $0.5285 \pm 0.0397$ \\
\quad Dermatofibroma (1.15\%)                 & $0.3359 \pm 0.0961$ \\
\quad Melanoma (11.11\%)                      & $0.4551 \pm 0.0349$ \\
\quad Melanocytic nevi (67.05\%)              & $0.8896 \pm 0.0088$ \\
\quad Vascular lesions (1.41\%)               & $0.6807 \pm 0.0680$ \\
\bottomrule
\end{tabular}
\end{table}

The per-class breakdown reveals the task's structure clearly. The
majority class \emph{melanocytic nevi} (67\% prevalence) is the
easiest to learn, reaching $F1 = 0.890$ with very small across-seed
variance ($\pm 0.009$). The two rarest classes,
\emph{dermatofibroma} ($1.15\%$) and \emph{vascular lesions}
($1.41\%$), show the largest across-seed variance ($\pm 0.096$ and
$\pm 0.068$), reflecting the difficulty of stably learning classes
represented by only $80$ and $99$ training samples respectively.
The two clinically critical classes \emph{melanoma} and
\emph{basal cell carcinoma} reach mean F1 of $0.455$ and $0.538$
respectively --- well above chance but substantially below the
majority class, consistent with the medical-imaging literature on
small CNNs at $28\times28$ resolution.

The macro-F1 of $0.56$ is a benchmark-typical result for
DermaMNIST at this resolution and architecture, comparable to
published values on the MedMNIST v2 benchmark. It is not intended
as a clinical-grade performance estimate; the input resolution is
deliberately compressed and the architecture is a small CNN, both
chosen for comparability across the federated experiments in this
dissertation rather than for clinical deployment.

\subsection{Role in the subsequent analysis}
\label{sec:centralised-role}

The centralised baseline serves three purposes throughout the rest
of the dissertation.

First, it is the \emph{upper bound} against which federated
configurations are measured. Subsequent results report recovery
percentages relative to the $0.5600$ centralised mean
(e.g.\ ``FedProx recovers $90.7\%$ of centralised performance''),
which makes the absolute magnitude of federation costs
interpretable. The within-federation algorithmic differences
discussed in Sections~\ref{sec:results-statistical-het}
and~\ref{sec:results-system-het} are then framed against this
fixed reference rather than against each other in isolation.

Second, the across-seed variance of the centralised baseline
($\pm 0.024$ macro-F1) sets a noise floor below which seed-level
differences in federated experiments cannot be reliably
distinguished from optimisation noise alone. Federated results
with within-pair $\Delta$ values smaller than this in magnitude
should be read as descriptive rather than substantive.

Third, the per-class profile in
Table~\ref{tab:centralised-results} establishes the
\emph{intrinsic difficulty} of each class on this task,
independent of any federation effect. The federated per-class
breakdowns reported later in the dissertation can be interpreted
relative to this profile: a class that already shows high
variance in the centralised setting is unlikely to support strong
inferential claims in the federated setting either.

\subsection{Summary}
\label{sec:centralised-summary}

The centralised baseline establishes that the $28\times28$
DermaMNIST classification task with the \texttt{DermMNISTCNN}
architecture supports a mean test macro-F1 of approximately $0.56$
under non-federated training, with the majority class melanocytic
nevi at near-ceiling performance ($F1 \approx 0.89$) and the
rarest classes at substantially lower and noisier per-class F1.
This serves as the reference against which all federated results
in the remainder of the dissertation are interpreted.

% ==============================================================================

\section{Centralised Baseline}
\label{sec:centralised-baseline}

Before evaluating any federated configuration, we establish a
centralised non-federated reference. The centralised baseline trains
the same \texttt{DermMNISTCNN} architecture
(Section~\ref{sec:model}) on the pooled DermaMNIST training set with
no federation. It provides an upper bound on what can be achieved on
this task with this architecture, against which the federated
results in Section~\ref{sec:fl-results} are interpreted as
fractional recoveries of the centralised performance.

\subsection{Training setup}
\label{sec:centralised-setup}

The centralised baseline pools all 7{,}007 training samples, the
1{,}003 validation samples, and the 2{,}005 test samples described
in Section~\ref{sec:dataset} into a single non-federated training
loop. The model architecture is held identical to the federated
experiments (Section~\ref{sec:model}): a four-block GroupNorm CNN
with 423{,}175 trainable parameters operating on $28\times28$ RGB
inputs. The optimiser is SGD with learning rate $0.01$, momentum
$0.9$, weight decay $0.0$, and batch size $32$ --- the same
hyperparameters used by each client in the federated experiments,
chosen for direct comparability rather than for centralised
optimum. The model is trained for $50$ epochs, with the test set
evaluated once at the epoch of highest validation macro-F1
(\emph{best-validation checkpoint} protocol).

To match the federated experiments and support a paired-seed
analysis where applicable, the centralised baseline is repeated
across the same ten random seeds used for the federated runs:
$\{42,\,123,\,456,\,789,\,999,\,2024,\,31337,\,8675309,\,161803,\,$
$271828\}$. Within each seed the data shuffle, model initialisation,
and training trajectory are deterministic; the seed-level variation
quantifies optimisation noise on this task.

\subsection{Results}
\label{sec:centralised-results}

Table~\ref{tab:centralised-results} reports test performance at the
best-validation checkpoint across the ten seeds. The centralised
model reaches a mean test macro-F1 of $0.560\pm0.024$ with mean
classification accuracy $0.764\pm0.012$ and mean balanced accuracy
$0.561\pm0.034$. The per-seed test macro-F1 ranges from $0.518$ to
$0.597$. The best-validation epoch is consistently in the second
half of training, with a median of $47$ out of $50$ epochs, so the
$50$-epoch training budget is appropriate for this task and
architecture.

\begin{table}[!htbp]
\centering
\caption{Centralised baseline test performance at the
best-validation checkpoint, reported as mean $\pm$ standard
deviation across $n=10$ seeds. Per-class F1 is the F1 score on
each of the seven DermaMNIST classes; class prevalence in the
training set is shown in parentheses.}
\label{tab:centralised-results}
\small
\begin{tabular}{lc}
\toprule
Metric & Test value (mean $\pm$ SD) \\
\midrule
Macro-F1                        & $\mathbf{0.5600 \pm 0.0240}$ \\
Balanced accuracy               & $0.5608 \pm 0.0341$ \\
Overall accuracy                & $0.7637 \pm 0.0123$ \\
Cross-entropy loss              & $0.7778 \pm 0.0653$ \\
\midrule
\multicolumn{2}{l}{\textit{Per-class F1:}} \\
\quad Actinic keratoses (3.27\%)              & $0.4923 \pm 0.0433$ \\
\quad Basal cell carcinoma (5.13\%)           & $0.5380 \pm 0.0263$ \\
\quad Benign keratosis (10.97\%)              & $0.5285 \pm 0.0397$ \\
\quad Dermatofibroma (1.15\%)                 & $0.3359 \pm 0.0961$ \\
\quad Melanoma (11.11\%)                      & $0.4551 \pm 0.0349$ \\
\quad Melanocytic nevi (67.05\%)              & $0.8896 \pm 0.0088$ \\
\quad Vascular lesions (1.41\%)               & $0.6807 \pm 0.0680$ \\
\bottomrule
\end{tabular}
\end{table}

The per-class breakdown reveals the task's structure clearly. The
majority class \emph{melanocytic nevi} (67\% prevalence) is the
easiest to learn, reaching $F1 = 0.890$ with very small across-seed
variance ($\pm 0.009$). The two rarest classes,
\emph{dermatofibroma} ($1.15\%$) and \emph{vascular lesions}
($1.41\%$), show the largest across-seed variance ($\pm 0.096$ and
$\pm 0.068$), reflecting the difficulty of stably learning classes
represented by only $80$ and $99$ training samples respectively.
The two clinically critical classes \emph{melanoma} and
\emph{basal cell carcinoma} reach mean F1 of $0.455$ and $0.538$
respectively --- well above chance but substantially below the
majority class, consistent with the medical-imaging literature on
small CNNs at $28\times28$ resolution.

The macro-F1 of $0.56$ is a benchmark-typical result for
DermaMNIST at this resolution and architecture, comparable to
published values on the MedMNIST v2 benchmark. It is not intended
as a clinical-grade performance estimate; the input resolution is
deliberately compressed and the architecture is a small CNN, both
chosen for comparability across the federated experiments in this
dissertation rather than for clinical deployment.

\subsection{Role in the subsequent analysis}
\label{sec:centralised-role}

The centralised baseline serves three purposes throughout the rest
of the dissertation.

First, it is the \emph{upper bound} against which federated
configurations are measured. Subsequent results report recovery
percentages relative to the $0.5600$ centralised mean
(e.g.\ ``FedProx recovers $90.7\%$ of centralised performance''),
which makes the absolute magnitude of federation costs
interpretable. The within-federation algorithmic differences
discussed in Sections~\ref{sec:results-statistical-het}
and~\ref{sec:results-system-het} are then framed against this
fixed reference rather than against each other in isolation.

Second, the across-seed variance of the centralised baseline
($\pm 0.024$ macro-F1) sets a noise floor below which seed-level
differences in federated experiments cannot be reliably
distinguished from optimisation noise alone. Federated results
with within-pair $\Delta$ values smaller than this in magnitude
should be read as descriptive rather than substantive.

Third, the per-class profile in
Table~\ref{tab:centralised-results} establishes the
\emph{intrinsic difficulty} of each class on this task,
independent of any federation effect. The federated per-class
breakdowns reported later in the dissertation can be interpreted
relative to this profile: a class that already shows high
variance in the centralised setting is unlikely to support strong
inferential claims in the federated setting either.

\subsection{Summary}
\label{sec:centralised-summary}

The centralised baseline establishes that the $28\times28$
DermaMNIST classification task with the \texttt{DermMNISTCNN}
architecture supports a mean test macro-F1 of approximately $0.56$
under non-federated training, with the majority class melanocytic
nevi at near-ceiling performance ($F1 \approx 0.89$) and the
rarest classes at substantially lower and noisier per-class F1.
This serves as the reference against which all federated results
in the remainder of the dissertation are interpreted.
